% ===========================================================================
\section{剩余类、欧拉函数与欧拉定理}\label{sec:euler}

\block{等价类与代表元}
模 $m$ 的同余关系把 $\Z$ 分成等价类：
\[
  \overline{0}=\{\ldots,-m,0,m,\ldots\},\quad
  \overline{1}=\{\ldots,-m+1,1,m+1,\ldots\},\quad\ldots\quad
  \overline{i}=\{i+km : k\in\Z\}.
\]
取每一类的最小非负代表元，得到\keyword{完全剩余系}
\[
  \Z_m=\{\overline{0},\overline{1},\ldots,\overline{m-1}\}.
\]
在剩余类上规定运算
\[
  \overline{a}+\overline{b}=\overline{a+b},\qquad
  \overline{a}\cdot\overline{b}=\overline{ab}.
\]
由同余与加乘的相容性（第 \ref{sec:cong} 节），运算结果不依赖代表元的选取，
即运算是良定义的。

\begin{definition}[既约剩余系与欧拉函数]
\[
  \Z_m^{*}=\{\,\overline{c}\in\Z_m : \gcd(c,m)=1\,\}
\]
称为模 $m$ 的\keyword{既约剩余系}；其中每个元素都可逆，且
$\overline{a}^{-1}=\overline{a^{-1}}$。其元素个数
\[
  \varphi(m)=|\Z_m^{*}|
\]
称为\keyword{欧拉函数}。
\end{definition}

\begin{theorem}[欧拉定理]\label{thm:euler}
对任意 $n,a$ 满足 $\gcd(n,a)=1$（即 $a\cop n$），有
\[
  a^{\varphi(n)}\equiv1\pmod n.
\]
\end{theorem}

\begin{proof}
因为 $a$ 在模 $n$ 下可逆，映射 $\overline{i}\mapsto\overline{i}\cdot\overline{a}$
是 $\Z_n^{*}$ 上的一个双射（逆映射为乘以 $\overline{a}^{-1}$），即“乘以 $a$”
只是把既约剩余系中的元素重新排列。于是
\[
  \prod_{i\in\Z_n^{*}} i \;\equiv\; \prod_{i\in\Z_n^{*}}(i\cdot a)
  \;\equiv\; a^{\varphi(n)}\prod_{i\in\Z_n^{*}} i \pmod n.
\]
乘积 $\prod i$ 仍属于 $\Z_n^{*}$，故可逆；两侧同乘其逆即得
\[
  1\equiv a^{\varphi(n)}\pmod n. \qedhere
\]
\end{proof}

\begin{corollary}[费马小定理]\label{cor:fermat}
对素数 $p$，若 $p\nmid a$，则
\[
  a^{p-1}\equiv1\pmod p.
\]
\end{corollary}

\begin{proof}
$p$ 为素数时 $\varphi(p)=p-1$（$1,\ldots,p-1$ 都与 $p$ 互素），
代入欧拉定理即得。
\end{proof}

\begin{corollary}\label{cor:fermat-all}
对素数 $p$ 与任意整数 $a$，
\[
  a^{p}\equiv a\pmod p.
\]
\end{corollary}

\begin{proof}
若 $p\nmid a$，在费马小定理两边同乘 $a$；若 $p\mid a$，两边都 $\equiv0$。
\end{proof}
